\subsection{A reconciled company case}

One hypothetical company is used throughout the chapter so that ratios, quality indicators, forecasts and decisions can be traced to the same statements. Its numbers are illustrative, not data about any real company.

\begin{qfexample}[Harbour Tools plc, 2024--2025: income statements, balance sheets, cash flows and equity movements reconciled from one set of transactions]\label{ex:harbour-tools-statements}
\qffield{Statement} \emph{The company.} Harbour Tools plc is a hypothetical European maker of power tools reporting under IFRS as adopted in the EU, in EUR millions, with calendar years. Before IFRS~18 (\hyperref[meth:accounting-change-effects]{2027 transition}), it presents expenses by function. It classifies interest paid as an operating cash flow, which IAS~7 then permitted \citep[IAS~7 paras.~31--33]{eu-2023-1803-ias7}. It has:
\begin{itemize}
\item a term loan at 4 percent repaid at 25 a year and a revolving credit facility (RCF) of 100;
\item leases accounted for under IFRS~16 at a 5~percent rate;
\item a funded defined benefit plan with a 4~percent discount rate;
\item equity-settled share options;
\item capitalised development costs and goodwill of 120 from an earlier acquisition.
\end{itemize}
The tax rate is 25~percent. Tax depreciation exceeds book depreciation; the share-based expense is not deductible; the 2025 restructuring provision is deductible when paid.

\begin{center}\small
\begin{tabular}{@{}lrr@{\qquad}lrr@{}}
\toprule
Income statement & 2025 & 2024 & Cash flows & 2025 & 2024\\
\midrule
Revenue & 1{,}080.00 & 1{,}000.00 & Received from customers & 1{,}065.00 & 985.00\\
Cost of sales & $-697.00$ & $-640.00$ & Paid to suppliers & $-479.00$ & $-442.00$\\
Gross profit & 383.00 & 360.00 & Paid to employees & $-283.00$ & $-268.00$\\
Selling and administration & $-233.00$ & $-218.00$ & Other S\&A, research, pension & $-104.00$ & $-99.00$\\
Research and development & $-38.00$ & $-35.00$ & Restructuring paid & $-6.00$ & 0.00\\
Restructuring & $-18.00$ & 0.00 & Interest paid & $-14.00$ & $-16.00$\\
Disposal gain (loss) & $-1.00$ & 3.00 & Income tax paid & $-16.96$ & $-15.50$\\
Operating profit & 93.00 & 110.00 & \textbf{Operating cash flow} & \textbf{162.04} & \textbf{144.50}\\
Finance costs & $-16.16$ & $-18.00$ & Capex, development, disposals & $-124.00$ & $-102.00$\\
Profit before tax & 76.84 & 92.00 & Loan, lease principal, RCF & $-10.00$ & $-45.00$\\
Income tax & $-21.46$ & $-25.00$ & Dividends; options exercised & $-26.00$ & $-22.00$\\
\textbf{Profit} & \textbf{55.38} & \textbf{67.00} & \textbf{Change in cash} & \textbf{2.04} & $\mathbf{-24.50}$\\
OCI: pension remeasurement & 4.00 & $-6.00$ & & &\\
\bottomrule
\end{tabular}
\end{center}

\begin{center}\small
\begin{tabular}{@{}lrrr@{\qquad}lrrr@{}}
\toprule
31 December & 2025 & 2024 & 2023 & & 2025 & 2024 & 2023\\
\midrule
Cash & 37.54 & 35.50 & 60.00 & Trade payables & 124.00 & 118.00 & 110.00\\
Trade receivables & 180.00 & 165.00 & 150.00 & Accrued wages; tax & 48.00 & 46.00 & 42.00\\
Inventories & 212.00 & 200.00 & 180.00 & Restructuring provision & 12.00 & 0.00 & 0.00\\
PP\&E, cost & 1{,}040.00 & 955.00 & 900.00 & Lease liabilities & 65.00 & 60.00 & 80.00\\
Accumulated depreciation & $-440.00$ & $-390.00$ & $-360.00$ & Term loan; RCF & 290.00 & 275.00 & 300.00\\
Right-of-use assets & 66.00 & 60.00 & 80.00 & Net pension deficit & 48.16 & 54.00 & 50.00\\
Development costs & 60.00 & 50.00 & 40.00 & Deferred tax liability & 41.00 & 37.50 & 30.00\\
Goodwill & 120.00 & 120.00 & 120.00 & Equity & 647.38 & 605.00 & 558.00\\
\textbf{Total assets} & \textbf{1{,}275.54} & \textbf{1{,}195.50} & \textbf{1{,}170.00} & \textbf{Total} & \textbf{1{,}275.54} & \textbf{1{,}195.50} & \textbf{1{,}170.00}\\
\bottomrule
\end{tabular}
\end{center}

\emph{Main transactions.}
\begin{itemize}
\item Cost of sales absorbs production depreciation (60 and 66) and production labour, as IAS~2 requires for conversion costs \citep[IAS~2 para.~12]{eu-2023-1803-ias2}.
\item Selling and administration include right-of-use depreciation (20 and 24), the share-based expense (8 and 9) and pension service cost (10 and 11).
\item Research and development comprise research expensed and amortisation of capitalised development (10 and 12).
\item Finance costs comprise loan interest, lease interest and net pension interest.
\item In 2025 a lease of 30 was added without cash, and 40 of the RCF was drawn on 31~December.
\end{itemize}

\emph{Equity movements, 2025.} Opening 605.00, plus profit 55.38, plus OCI 4.00, plus share-based expense 9.00, plus options exercised for cash 6.00, less dividends 32.00, gives 647.38. Retained earnings move from 400.00 to 423.38, which is 400.00 plus 55.38 less 32.00.

\qffield{Notation} EBITDA here is operating profit plus depreciation of PP\&E and right-of-use assets plus amortisation: 195 in 2025 (200 in 2024). Debt is the term loan, RCF and lease liabilities: 355 at the end of 2025.

\qffield{Conditions} Hypothetical; built to illustrate mechanics and analysis, not representative of an industry. Tax deductibility is simplified: pension cost is deductible as recognised, the share-based expense never. Current portions are taken as next year's scheduled repayments (term loan 25; lease principal 24.75). The RCF matures beyond 2028 with a right to roll over drawings, so it is non-current \citep[Annex, IAS~1 para.~73]{eu-2023-2822-ias1}.

\qffield{Practice} \emph{Independent reconciliations.} The statements are built as a double-entry ledger. Four checks then pass for both years:
\begin{itemize}
\item the direct-method operating cash flow (sum of cash postings) equals the indirect-method figure derived only from the income statement and balance-sheet changes (162.04 in 2025);
\item operating, investing and financing cash flows sum to the change in cash;
\item every equity component rolls forward;
\item the balance sheet balances.
\end{itemize}
The 2025 indirect reconciliation is: profit 55.38; plus non-cash charges 114.50 (depreciation 66, right-of-use depreciation 24, amortisation 12, share-based expense 9, deferred tax 3.50); plus the disposal loss 1.00; less pension contributions above P\&L cost $-1.84$; less working capital and provision changes $-7.00$. The total is 162.04.

\emph{First reading.}
\begin{itemize}
\item Revenue grew 8 percent, but operating profit fell from 110 to 93 because of an 18 restructuring charge and a 0.54-point lower gross margin.
\item Operating cash flow rose, partly because 12 of the restructuring charge was still unpaid.
\item Capital spending of 127 (net of disposals 124) consumed most of operating cash flow, so dividends of 32 were funded partly by the year-end RCF drawing.
\end{itemize}
The ratio, quality and forecast entries develop each of these points.

\qffield{Pitfalls} Reading the 2025 operating cash flow as a sustainable improvement without the 12 unpaid provision. Treating the year-end RCF drawing as a sign of distress without the capital spending and dividends that caused it.

\qffield{See also} \hyperref[meth:statement-linked-cash-flows]{linking the statements}; \hyperref[meth:income-statement-ratios]{income-statement ratios}; \hyperref[ex:harbour-pro-forma-model]{forecast}.

\end{qfexample}
